\documentclass[11pt]{article}
\usepackage[utf8]{inputenc}
\usepackage{amsmath,amssymb,amsfonts,amsthm}
\usepackage{graphicx}
\usepackage{xcolor}
\usepackage{longtable}
\usepackage{hyperref}
\usepackage{tikz}
\usetikzlibrary{calc, arrows.meta, decorations.markings, patterns, shapes.geometric, positioning, 3d, decorations.pathmorphing}
\usepackage{geometry}
\usepackage{booktabs}
\PassOptionsToPackage{hyphens}{url}
\usepackage{tabularx}
\geometry{margin=1in}
\emergencystretch=3em
\Urlmuskip=0mu plus 1mu\relax
\def\UrlBreaks{\do\/\do\-\do\.\do\_}
\newenvironment{acknowledgments}
  {\section*{Acknowledgments}}
  {}
\usepackage{cite}

% Theorem environments
\newtheorem{theorem}{Theorem}
\newtheorem{lemma}{Lemma}
\newtheorem{proposition}{Proposition}
\newtheorem{corollary}{Corollary}
\newtheorem{definition}{Definition}
\newtheorem{remark}{Remark}
\newtheorem{postulate}{Postulate}
\newtheorem{assumption}{Assumption}

\providecommand{\KL}{\operatorname{KL}}
\providecommand{\softmax}{\operatorname{softmax}}

\begin{document}

\title{A Gauge-Theoretic Framework Toward a Participatory ``It From Bit'' Program}

\author{
  Robert C. Dennis\thanks{Corresponding author. Email: cdenn016@gmail.com}\\[0.3cm]
  \small Independent Researcher\\
  \small Leander, Texas, United States\\[0.2cm]
  \small ORCID: \href{https://orcid.org/0009-0004-4688-2341}{0009-0004-4688-2341}
}

\date{\today}

\maketitle

\begin{abstract}
Wheeler proposed, as a working hypothesis, that every physical quantity derives its significance from registered yes-or-no answers in a universe whose observers participate in what they observe, and he asked how one world arises out of many observer-participants. This paper gives a classical, multi-agent, information-geometric model of four parts of that program. Agents are local sections of a principal bundle at one context point; each holds a probability law in its own frame, and agents compare laws only through explicit transports whose residual relative entropy is the frame-independent part of their disagreement. First, registered change as duration: each agent's clock is the Fisher length of its belief trajectory. Second, observer-dependent geometry: metrics on an index manifold without a physical metric are pulled back from belief sections, a step that uses base derivatives and is developed in a companion supplement. Third, one world out of many: agreement is a stationary point of a consensus free-energy functional whose attention weights are Gibbs distributions over transported divergences. Fourth, participation across scales: exact coarse-graining assembles coherent clusters into meta-agents whose laws act as lagged priors for their members, and the hierarchy closes on a declared root, a fixed top-level prior. A claim-status table separates theorems, derivations, structural correspondences, and conjectures. The model has no irreversible registration and no complementarity, fixes its law in advance, presupposes a continuum and an ordering parameter, and relaxes unless driven by observations; how come the quantum lies outside its scope.
\end{abstract}

\noindent\textbf{Keywords:} Participatory realism $\cdot$ It from bit $\cdot$ Gauge theory $\cdot$ Information geometry $\cdot$ Variational free energy $\cdot$ Multi-agent inference

\tableofcontents
\clearpage

\part*{Part I: The Construction}\addcontentsline{toc}{part}{Part I: The Construction}

\section{Introduction}\label{sec:introduction}

In an essay delivered in 1989 and published in 1990, John Archibald Wheeler returned to ``the age-old question, How come existence?'' and offered as a working hypothesis that every physical quantity, every ``it'', derives its significance from ``binary yes-or-no indications'' elicited by observing devices, in a universe in which observers take part in bringing about what they observe~\cite{Wheeler1990, Wheeler1983}. Among the questions he posed was ``How come `one world' out of many observer-participants?'', and he adopted F{\o}llesdal's dictum that meaning is the joint product of all the evidence available to those who communicate~\cite{Wheeler1990}. In physics the program has been pursued chiefly as a reconstruction of quantum theory from informational principles. Wheeler himself credited Wootters, whose statistical distance between preparations, defined by the distinguishability of measurement outcomes, coincides up to a constant factor with the Fisher--Rao distance~\cite{Wootters1981}. Zeilinger took as foundational that an elementary system carries one bit~\cite{Zeilinger1999}; Hardy, Chiribella, D'Ariano and Perinotti, and Masanes and M{\"u}ller derived the formalism of quantum theory from operational and informational axioms~\cite{Hardy2001, ChiribellaDArianoPerinotti2011, MasanesMueller2011}; H{\"o}hn reconstructed qubit quantum theory from an observer who interrogates a system with binary questions~\cite{Hoehn2017}; and M{\"u}ller derived the appearance of an external world with simple probabilistic laws from observer states alone, under the title ``Law without law''~\cite{Mueller2020}. QBism claims Wheeler's participatory program by name as participatory realism~\cite{Fuchs2017}, relational quantum mechanics relativizes states to arbitrary physical systems~\cite{Rovelli1996}, and Caticha's entropic program treats space as a statistical manifold and quantum dynamics as entropic inference~\cite{Caticha2005, Caticha2019spacetime, Caticha2019}. The present framework is not a reconstruction of quantum theory. It is classical and multi-agent, and it supplies two objects those programs do not address: an explicit law of transport between the reference frames of different agents, and a consensus functional whose stationary points model agreement among them.

Two further lines bear on the construction. The Kant--Helmholtz line reads perception as structured by the perceiver and as unconscious statistical inference~\cite{Kant1781, Helmholtz1867}, and its modern descendant, the free-energy principle, derives perception, learning, and action from one variational free-energy functional~\cite{Friston2010, Friston2017, parr2022active}. For a physicist, a variational free energy is the negative evidence lower bound of a Bayesian model, the Gibbs--Bogoliubov functional minimized in a mean-field approximation, and the framework uses the term only in that technical sense; the consensus functional of Section~\ref{sec:free_energy} has the energy-minus-entropy form of a free energy but is a declared composite potential, not such a bound. Independently, spacetime has been argued to be emergent from thermodynamic, holographic, and entanglement structure~\cite{Jacobson1995, Padmanabhan2010, Bousso2002, Maldacena1999, VanRaamsdonk2010, MaldacenaSusskind2013, Swingle2012}. The framework sits between these readings. Agents with different beliefs pull back different geometries from one shared index space, a presupposed smooth index manifold $\mathcal{C}$ of contexts carrying no physical metric, called noumenal by analogy only (Section~\ref{sec:one_rung}), and no agent-independent geometry is singled out. The commitment concerns the structure of the map from this index space to what each agent represents; it does not stipulate that anything beyond $\mathcal{C}$ is inaccessible.

The organizing claim of the paper is a correspondence with Wheeler's program, stated together with its limits. Agents are local sections of a principal $G$-bundle at one base point $c_*$ of $\mathcal{C}$, each carrying a belief $q_i$, a prior $p_i$, a frame $\phi_i \in \mathfrak{gl}(K_q)$, and a slow generative model $s_i$ with hyper-prior $r_i$ (Section~\ref{sec:one_rung}); they are coupled through the transport $\Omega_{ij} = \exp(\phi_i)\exp(-\phi_j)$, which acts on covariances by $\Sigma \mapsto \Omega_{ij}\Sigma\Omega_{ij}^\top$, and through the consensus free-energy functional (Section~\ref{sec:free_energy}). On this construction the paper gives a classical, multi-agent, information-geometric model of four parts of Wheeler's program, labeled W1--W4 in Table~\ref{tab:wheeler_correspondence}. Observer-relative description, the comparison of laws only through an explicit transport whose residual divergence is the frame-independent part of a disagreement, underlies all four and is what makes the geometry of W2 observer-dependent. Derived duration and geometry are W1 and W2: the Fisher length of each agent's belief trajectory, and the metrics pulled back to $\mathcal{C}$ from belief sections, a step that uses base derivatives and so leaves the single base point (Section~\ref{sec:it_from_bit}). One world out of many observer-participants is W3: intersubjective agreement as a stationary point of the consensus functional, a formal counterpart of Wheeler's use of F{\o}llesdal's dictum. Participation across scales, W4, is modeled by one mechanism: an exact coarse-graining assembles a coherent cluster of level-$\ell$ agents into a level-$(\ell+1)$ meta-agent, whose law returns to its members as a lagged predictive prior, and the hierarchy is closed by a declared root, a fixed top-level prior acting as its boundary condition, together with an energy-level apex for the belief channel that pools current beliefs and is not a generative factor (Sections~\ref{sec:meta_agent_formation}--\ref{sec:ouroboros}). The coarse-graining is an exact block-variable pushforward in Kadanoff's sense~\cite{kadanoff1966scaling, Kadanoff1975, KadanoffHoughtonYalabik1976} followed by a truncation. It is not a renormalization-group transformation: it involves no rescaling, no flow on a fixed space of couplings, and no fixed point, and on a finite network no typed comparison of interaction spaces across a reducing step exists~\cite[Ch.~11]{DennisGaugeVFE2026}. Table~\ref{tab:wheeler_correspondence} lists the five parts W1--W5 of the correspondence with their status, and Table~\ref{tab:claim_status} records the status of every load-bearing claim.

The correspondence keeps four of Wheeler's commitments and drops four. It keeps the observer-relativity of description, the formation of one shared world out of many participants, the derivation of metric and duration from informational structure, and the participation of collectives in the priors of their members. It drops, first, the elementary act of registration. Beliefs here are continuous Gaussian laws, and the relative entropy of an update, measured in bits, is a unit of measure that is not additive under refinement of the update path. Section~\ref{sec:time_info} adds a binary registration channel whose expected information gain is at most one bit per registration and adds over conditionally independent registrations, which gives the word bit its physical meaning; even so, the construction contains no irreversible amplification and no complementarity, and its agents receive observations without choosing which question to ask, so ``no question, no answer'' is not represented. It drops, second, ``no continuum'': the index manifold $\mathcal{C}$ and the fibers are smooth and presupposed, and the regress that Wheeler answered with a loop is stopped here by an index floor and a declared root, a deliberate departure whose formal residue is discussed in Section~\ref{sec:ouroboros}. It drops, third, ``law without law'': the functional, the belief family, and the divergence are fixed in advance, and M{\"u}ller's observer-first program~\cite{Mueller2020} is closer to Wheeler on this point. It drops, fourth, the self-excited circuit. The autonomous dynamics relax to stationary points, observations supply the only persistent drive, and the geometry induced on $\mathcal{C}$ enters no term of the free energy, so the loop from information back to the participants closes only in the space of beliefs; closing it through the induced geometry requires base transport and is deferred to the companion supplement~\cite{DennisGaugeVFESuppB2026}. How come the quantum is outside the scope of the construction.

\begin{table}[t]
\centering
\small
\caption{The correspondence between strands of Wheeler's participatory program and the construction. Status codes are those of Section~\ref{sec:epistemic_status}; ``not claimed'' marks what the construction does not deliver.}
\label{tab:wheeler_correspondence}
\begin{tabularx}{\textwidth}{@{}>{\raggedright\arraybackslash}p{0.6cm}>{\raggedright\arraybackslash}p{2.4cm}>{\raggedright\arraybackslash}X>{\raggedright\arraybackslash}p{2.0cm}>{\raggedright\arraybackslash}p{2.6cm}@{}}
\toprule
 & Wheeler's strand & Formal counterpart here & Where & Status \\
\midrule
W1 & Registered change as the measure of duration & Fisher-length clock of each agent; exact clock--dissipation identity along an autonomous single-rate natural-gradient flow; binary registration channel whose expected information gain is at most one bit per registration; identification of the clock with time is postulate P1 & Section~\ref{sec:time_info} & P (identity), D (registration), S (reading as time) \\
W2 & Observer-dependent ``it'' & Pullback of the Fisher metric along belief sections; the covariant Fisher block is gauge invariant, so differences between agents come from differences of belief & Section~\ref{sec:it_from_bit} & D (pullback), P (invariance), S (reading as space) \\
W3 & One world out of many observer-participants & Stationary points of the consensus free energy; Gibbs attention rows with reduced value $-\tau\log Z$; consensus geometry, whose orbit average needs a regulator only on the frame-twist block & Sections~\ref{sec:free_energy} and~\ref{sec:collective_geometry} & P (rows, regulator obstruction), S (reading as one world), C (regulated metric) \\
W4 & Participation across scales; no tower of turtles & Exact coarse-graining into meta-agents; lagged predictive priors from parents to members; a declared root closing the hierarchy & Sections~\ref{sec:meta_agent_formation}--\ref{sec:ouroboros} & P (coarse-graining, root), D (lagged prior), C (formation by the agents' own inference) \\
W5 & What is not delivered & Discrete irreversible registration with complementarity; law without law; no continuum; a self-excited circuit; how come the quantum & Sections~\ref{sec:scope_limitations} and~\ref{sec:open_problems} & not claimed \\
\bottomrule
\end{tabularx}
\end{table}

\paragraph{Interpretive stance.}\label{sec:cognitive_first} The author's interpretive stance carries status S throughout, and no equation of the paper depends on it. It is a cognitive-first structural monism: the primitives of the construction are the agents' beliefs and generative models, and physics-like structure appears as a derived pullback over a non-cognitive index space $\mathcal{C}$, which plays the structural role of Whitehead's extensive continuum~\cite{Whitehead1929}. The stance inverts the Primacy of Physics Constraint of Ladyman and Ross's naturalized metaphysics~\cite{Ladyman2007} while sharing the relational emphasis of ontic structural realism~\cite{Ladyman1998} (Section~\ref{sec:discussion}). It makes four commitments. Every system at every level carries frame structure and takes part in variational dynamics on a statistical manifold suited to its level: a multivariate Gaussian family on $\mathbb{R}^{K_q}$ in the classical construction developed here, and, at level zero, the manifold of density operators with the Umegaki quantum relative entropy as divergence, whose construction is open (Section~\ref{sec:open_problems}). All levels share the transport structure of a principal bundle with a structure group suited to the level~\cite{Nakahara2003}, $\mathrm{GL}(K_q,\mathbb{R})$ in the classical construction and $U(\mathcal{H})$ at level zero. The base $\mathcal{C}$ is an index manifold of contexts, not a physical thing. Agents at level zero are ascribed frame structure and variational dynamics, not phenomenal properties, and the combination problem this incurs is owned as open (Section~\ref{sec:consciousness}). Two remarks fix the status of these commitments. First, the commitments to universal agency and universal transport exclude no system: any dynamics can be given a trivial bundle, trivial frames, and a residual functional of which it is the minimizer, so, as with Kretschmann's objection to general covariance~\cite{Kretschmann1917, Norton1993}, they carry no empirical content by themselves. Second, at the level of kernels the commitment to universal agency is free: observations can always be presented as messages from boundary agents without changing any law, although that presentation removes no conditioning on the record~\cite[Ch.~5]{DennisGaugeVFE2026}; what the stance adds, a criterion of agency for every node, is not supplied by the mathematics. On this reading an electron would be an agent at level zero, a molecule a meta-agent assembled from atomic constituents, and a brain a meta-agent at a higher level; against the panpsychism taxonomy the position is panprotopsychist and constitutive (Section~\ref{sec:consciousness}). The vocabularies of Wheeler, Kant, and Lahav and Neemeh are borrowed with the departures listed above and in Section~\ref{sec:scope_limitations}. Symbol conventions, including the few deliberate overloadings, are collected in Appendix~\ref{sec:notation}.

\section{Claim Status and Scope}\label{sec:epistemic_status}

Every load-bearing claim carries one of five status codes, defined once here. P (proven) marks a claim proved in this paper or entailed by a cited result. D (derived) marks a result derived here under stated assumptions, such as a Gaussian family, a fixed schedule, or a fixed attention row. S (structural) marks a structural correspondence or an interpretation: a mathematical object is read as the formal counterpart of a physical or philosophical notion, and the reading adds no theorem. E (empirical) marks a measured result; no claim of the present paper has this status, and its multi-scale claims are tested by the companion's pairing-selection conjecture and by the prospective protocol of Section~\ref{sec:participatory}, its finite-sample counterpart. C (conjectured) marks a conjecture or an open problem. The numbered inputs P1--P10 of Part II are postulates, registered in Section~\ref{sec:part2_preamble}; a postulate is adopted, not derived, and a result built on one inherits its status. Part I declares one further postulate outside this register, the kinetic metric of Section~\ref{sec:velocity_quadratic}; no other result uses it, and the inertial reading of precision it expresses has status S. The letter P as a status code and the labels P1--P10 are distinct. Part I is the construction: its results carry status P or D, and its interpretive readings are marked S. Part II is the speculative program: its physical readings carry status S or C and rest on the postulates, while the mathematical results it uses (the clock--dissipation identity, the registration proposition, the pullback and its invariance, the signature constructions) carry status P or D and do not depend on those readings. Table~\ref{tab:claim_status} gives each load-bearing claim with its location, its status, and what would change that status. The body states a claim's status at its first statement and otherwise refers to the table.

{\small
\setlength{\LTcapwidth}{\textwidth}
\begin{longtable}{@{}>{\raggedright\arraybackslash}p{0.39\textwidth}>{\raggedright\arraybackslash}p{0.17\textwidth}>{\raggedright\arraybackslash}p{0.12\textwidth}>{\raggedright\arraybackslash}p{0.22\textwidth}@{}}
\caption{Status of the load-bearing claims. ``Theory'' denotes the companion manuscript~\cite{DennisGaugeVFE2026} and ``Supp.~B'' its base-transport supplement~\cite{DennisGaugeVFESuppB2026}.}\label{tab:claim_status}\\
\toprule
Claim & Where & Status & What would change it \\
\midrule
\endfirsthead
\toprule
Claim & Where & Status & What would change it \\
\midrule
\endhead
\bottomrule
\endfoot
\multicolumn{4}{@{}l}{\emph{Part I: the construction}}\\
$D_{\mathrm{KL}}(q_i\|\Omega_{ij}q_j)$ is invariant under a common $\mathrm{GL}(K_q)$ pushforward; in Regime~I every frame can be removed from the alignment, attention, and cross-scale terms at one base point & Theorem~\ref{thm:glk_invariance}; Proposition~\ref{prop:regime1_removability} & P & none (theorem) \\
At one base point the active symmetry is one global copy of $G$; per-agent frames are passive re-expressions & Section~\ref{sec:gauge_covariance_principle}; Theory Ch.~2 & P & none \\
Regime~I cocycle closure and the lattice Bianchi identity with orientation reversal, read as parallels to ``the boundary of a boundary is zero'' & Section~\ref{sec:boundary_of_boundary} & P (identities), S (parallel) & a conservation law derived from them \\
The attention row, a softmax over transported divergences, is the Gibbs optimum of the entropy-regularized consensus problem, with reduced value $-\tau\log Z$ and the envelope gradient & Section~\ref{sec:free_energy} & P & none \\
On an open mean-field family, live pairwise divergences cannot be the negative evidence bound of one fixed state-level joint; the consensus functional is a declared composite potential & Section~\ref{sec:free_energy}; Appendix~\ref{app:augmented_joint}; Theory Chs.~6 and~12 & P & none; the functional stays a declared ansatz \\
Among real-analytic $f$-divergences, the relative entropy with the receiver law in the first slot is the unique one giving the geometric-mean best response to frozen neighbors & Section~\ref{sec:mixture_derivation}; Appendix~\ref{app:conditional_representation} & P (best response only) & extension to the full functional fails \\
Observations as couplings to environmental agents & Section~\ref{sec:precision_envelope}; Theory Ch.~5 & D (Gaussian, full-rank linear readout); P (kernel level) & none \\
The functional is invariant under the change of reference $U_i\mapsto U_i g$ and, without observations, under the global belief rotation; a fixed observation kernel explicitly breaks the global belief-rotation symmetry & Section~\ref{sec:symmetry_breaking} & P (invariances), D (breaking) & none \\
Adaptive precision as a declared penalty & Section~\ref{sec:precision_envelope} & D & a normalized hyperprior \\
Adiabatic separation of fast beliefs and slow models & Section~\ref{sec:timescales} & C & verified Tikhonov--Fenichel hypotheses \\
The precision Hessian is a stiffness and the Fisher metric the friction tensor of the first-order natural-gradient flow & Section~\ref{sec:stiffness_friction} & D & none \\
Kinetic-metric postulate $M_K = M$, with its own second-order equation of motion, under which every mode has unit frequency & Section~\ref{sec:velocity_quadratic} & S (declared postulate, not among P1--P10) & a second-order dynamics derived from the functional \\
Moment-projection barycenter: closed form and M-step reading & Section~\ref{sec:meta_agent_formation} & P & none \\
The frame Karcher mean on $\mathrm{SO}(K_q)$ is unique for rotation angles below $\pi/2$ & Section~\ref{sec:meta_agent_formation} & P & none \\
The parent frame is not identifiable, and no left-equivariant permutation-symmetric map $\mathrm{GL}^+(K_q)^n\to\mathrm{GL}^+(K_q)$ exists & Section~\ref{sec:meta_agent_formation} & P & none \\
Licensing criterion, degenerate under exact parents (Theory Ch.~9); threshold detector; meta-agents read as belief-bearing agents & Sections~\ref{sec:meta_agent_variational} and~\ref{sec:meta_agent_threshold} & C (criterion), S (detector, reading) & a derivation of the detector from the criterion \\
Contraction of the meta-agent weight map at rate $O(\varepsilon)$ & Lemma~\ref{lem:weight_contraction} & D (sketch) & a complete proof \\
The cross-scale shadow is the one-episode-lagged predictive prior of the episode-indexed joint under a transported-copy transition; the same-time form is an energy-level simplification and not a message of the augmented joint & Section~\ref{sec:cross_scale_shadows}; Theory Ch.~12 & D (lagged), S (same-time) & none \\
Meta-agent pullback geometry is dominated by the constituents' wherever the parent frame introduces no horizontal defect & Section~\ref{sec:cross_scale_shadows}; Supp.~B Ch.~B2 & P (conditional) & none; base-direction work deferred \\
Mean-field solution of the augmented Gaussian joint & Appendix~\ref{app:augmented_joint} & P & a proof that it governs the shadow dynamics \\
Exact coarse-graining: a block-variable pushforward and a truncation, with no rescaling, flow, or fixed point & Section~\ref{sec:rg_status}; Appendix~\ref{sec:meta_agent_rg}; Theory Chs.~9 and~11 & P & a comparison map between coupling spaces \\
A declared root closes the hierarchy, and the apex is an energy-level closure; the root read as the formal residue of ``no tower of turtles'' & Section~\ref{sec:ouroboros}; Theory Ch.~12 & P (root), S (reading) & none \\
The apex loop has the feedback structure of Wheeler's self-excited circuit, without its self-excitation (W5) & Section~\ref{sec:ouroboros} & S & sustained non-stationary dynamics without exogenous input \\
Observation injection sustains non-stationarity; without it the flow relaxes & Section~\ref{sec:ouroboros} & C (sustained), P (relaxation) & a proof that observation injection sustains non-stationarity, or a counterexample \\
Collectives are constituted by their members' inference: pairing selection by slow learning of the attention prior & Section~\ref{sec:participatory}; Theory Ch.~14 & C & proof or counterexample \\
Interpretive stance (cognitive-first structural monism) & Section~\ref{sec:introduction} & S & none; not a result \\
\multicolumn{4}{@{}l}{\emph{Part II: the speculative program}}\\
Clock--dissipation identity along an autonomous single-rate natural-gradient flow & Section~\ref{sec:time_info}; Theory Ch.~7 & P & none; a driven flow adds $\partial_t\mathcal{F}$ (Eq.~\eqref{eq:clock_driven}) \\
Binary registration: $\mathbb{E}_O[D_{\mathrm{KL}}(p(x|O)\|p(x))] = I(X;O)\le H(O)\le 1$ bit, additive over conditionally independent registrations & Section~\ref{sec:time_info} & D & none \\
The information clock as the agent's time; the ordering parameter is presupposed & Section~\ref{sec:time_info} & S (postulate P1) & a calibration against a physical clock \\
Meta-clock ratio in the relative-entropy currency; a meta-agent accumulates no more Fisher duration than the collective of its constituents; the arc-length ratio & Section~\ref{sec:time_info}; Theory Ch.~7 & D (ratio), P (contraction), C (arc length) & a proof of the arc-length ratio \\
The Fisher pullback tensor and its rank; gauge invariance of the covariant Fisher block & Section~\ref{sec:it_from_bit}; Supp.~B Ch.~B2 & D (pullback), P (invariance) & none \\
The pullback read as a formal counterpart of ``it from bit'' & Section~\ref{sec:it_from_bit} & S & the topology or dimension of $\mathcal{C}$ derived from registrations \\
All-real frame-twist form of signature $(-,+,+,+)$, exhibited under P4; in Regime~I a property of a declared reference section with no gauge-invariant content & Section~\ref{sec:signature_resolution} & P (under P4); C (selection) & a dynamical selection mechanism \\
Orbit-averaged consensus metric; its regulator obstruction concerns the frame-twist block and the continuum and noncompact limits & Section~\ref{sec:collective_geometry}; Proposition~\ref{prop:regulator_obstruction} & P (obstruction), C (regulated metric) & a constructed regulator \\
Physical correspondences of Table~\ref{tab:physical_correspondences}; gravity as curvature of an induced geometry, a placeholder & Sections~\ref{sec:phenomenological_interpretation} and~\ref{sec:speculative_physics} & S & a field equation with a source \\
Consciousness as frame-dependent presentation & Section~\ref{sec:consciousness} & S, C & none; interpretive \\
Open problems C1--C5 are mathematical problems and C6 a derivation target; the physical readings remain interpretive & Section~\ref{sec:open_problems} & C & solutions of C1--C6 \\
\multicolumn{4}{@{}l}{\emph{Appendices}}\\
High-precision agent (the rock): the mean back-action is suppressed, while the stationary precision is reduced by a factor tending to $1/(1+\beta)$; the rock read as an observed object & Appendix~\ref{app:examples} & D; S (reading) & none \\
Precision-linear covariance fixed point: exists if and only if $\rho(L)<2$, and is then globally stable under the receiver-only flow; its reading as a population equilibrium reproducing the shared prior reflects the doubled entropy weight & Appendix~\ref{app:covariance_dynamics} & P; S (reading) & none \\
The frozen mean map is a contraction for isotropic couplings & Appendix~\ref{app:augmented_joint} & P (isotropic) & fails for anisotropic couplings \\
Nested state--configuration bound, the two-level case of the Theory's multiscale bound; the one-particle-irreducible and two-particle-irreducible fluctuation program & Appendix~\ref{sec:beyond_mean_field_gauge_vfe}; Theory Ch.~12 & P (bound), C (program) & a constructed effective action \\
Gibbs lift of a normalizable consensus energy to belief configurations & Appendix~\ref{sec:meta_agent_rg} & P & none \\
Covariance of the coarse-graining map under one global $g$ at one base point, extending to left-invariant metrics; mean-sector Schur lemma for mean-only linear couplings & Appendix~\ref{app:rigorous_rg} & P & none \\
One-loop shift $-\tfrac12 T_{\mathrm{cfg}}(|I|-1)\log\det\Sigma_I$ of the parent's entropy weight; force residual measured against the zero-temperature profile, a finite-temperature version failing & Appendix~\ref{sec:meta_agent_rg} & D & none \\
Schematic residual bound for the truncation & Appendix~\ref{app:rigorous_rg} & C & pinned constants and tightness \\
Detector retention, conditional on a positive saving term $A_I$ whose source is not identified; the body detector is not covered & Appendix~\ref{app:rigorous_rg} & P (conditional) & a source for $A_I$ \\
\end{longtable}
}

\subsection{Scope and Limitations}\label{sec:scope_limitations}

The claims of the paper are delimited as follows. The construction supplies formal counterparts, in the sense of Table~\ref{tab:wheeler_correspondence}, for strands of Wheeler's participatory universe, of Kant's distinction between appearance and what appears, of observer-dependent description, and of the Lahav--Neemeh proposal that consciousness is frame-relative. We do not derive quantum mechanics or general relativity. The signature construction of Section~\ref{sec:signature_resolution} exhibits a Lorentzian frame-twist form under postulate P4 but selects nothing, and the orbit-averaged consensus metric of Section~\ref{sec:collective_geometry} is not the only gauge-invariant geometric candidate: the Fisher block of the pullback, computed with each agent's own covariant derivative, is gauge invariant without any orbit average~\cite[Ch.~B2]{DennisGaugeVFESuppB2026}, so the orbit average and its regulator are needed only for the frame-twist block.

Earlier simulations of the meta-agent hierarchy are not used as evidence. Section~\ref{sec:participatory} states a prospective test, which requires an active slow model channel and a control arm in which the attention prior is frozen. The control is needed because, on the four-agent cycle of the companion theory, a collective forms from every Gaussian initial law whether or not the attention prior learns, the random initial source assignment seeding the pair~\cite[Ch.~14]{DennisGaugeVFE2026}; the occurrence of formation therefore does not by itself distinguish participation from seeding. Whether the discrete detector of Section~\ref{sec:meta_agent_threshold} and the continuous variational criterion, from which the detector is not derived, produce the same organization remains open.

Several further limitations each receive one statement. The framework yields no quantitative physical predictions: it produces no values of physical constants, masses, couplings, or dimensionless ratios, and its induced tensors are dimensionless, so a map from information-geometric quantities to physical units is missing. There is no quantum extension: the framework contains no Hilbert space, no Born rule, and no measurement postulate, so its parallel with QBism and relational quantum mechanics is one of stance (Section~\ref{sec:related_work}), and the density-operator construction at level zero is open (Section~\ref{sec:open_problems}). The objectivity the framework delivers is agreement between the pullback geometries of agents that share a belief family, not access to an agent-independent geometry (Section~\ref{sec:collective_geometry}). The discussion of consciousness relocates the hard problem without solving or dissolving it (Section~\ref{sec:consciousness}).

Testability is tiered. Mathematical claims internal to the framework are settled by proof or counterexample. The participatory reading of the formalism, that collectives are constituted by their members' inference and are not supplied, has one decidable test in the companion theory: a conjecture that slow learning of the attention prior selects which pair of agents on a homogeneous four-agent cycle forms a collective, at the level of the law, against a frozen-prior control whose pairing law stays symmetric~\cite[Ch.~14]{DennisGaugeVFE2026}. The conjecture is settled by proof or counterexample, not by observation, and its failure would reduce that reading to hierarchical modeling in that family. The participatory-physics reading of Part~II has no identified falsifying observation and may be unfalsifiable as currently stated (Section~\ref{sec:gauge_invariance_consensus}). The open problems are collected in Section~\ref{sec:open_problems}.

\subsection{Relation to the Companion Theory}\label{sec:companion_theory}

A companion manuscript~\cite{DennisGaugeVFE2026} develops the single-base-point core of this program with normalized laws throughout: one collective evidence bound for the agent population~\cite[Ch.~5]{DennisGaugeVFE2026}, attention as inference over a declared source label~\cite[Ch.~6]{DennisGaugeVFE2026}, Fisher-length information clocks~\cite[Ch.~7]{DennisGaugeVFE2026}, exact Markov coarse-graining of the agent network into meta-agents~\cite[Chs.~9 and~11]{DennisGaugeVFE2026}, the multiscale free energy with its episode-indexed joint and declared root~\cite[Ch.~12]{DennisGaugeVFE2026}, and formation theorems for declared families~\cite[Ch.~13]{DennisGaugeVFE2026}. Its supplements treat the multivariate-Gaussian case~\cite{DennisGaugeVFESuppG2026} and base transport, including the connection-relative pullback geometry behind Section~\ref{sec:it_from_bit}~\cite{DennisGaugeVFESuppB2026}. Where the companion supersedes a construction used here, namely the consensus functional, the barycenter and detector, the same-time cross-scale shadows, and the self-closing apex, we say so at the construction. The companion declines every physical identification made in Part~II~\cite[Chs.~1 and~14]{DennisGaugeVFE2026}: it identifies no Fisher duration with physical time, gives the context base no spacetime or Lorentzian structure, treats calling a boundary node an agent as a naming decision, and does not read the free energy as a physical action. Declining a claim is not refuting it. Part~II adopts these identifications as its postulates, so the two documents agree on the mathematics and differ only in what is postulated here. The companion's vocabulary translates as follows: a recognition law is a variational posterior, the source prior is the attention prior $\pi_{ij}$, a declared root is a top-level prior fixed as a boundary condition, and a Fisher duration is the information clock.

\section{Related Work}\label{sec:related_work}

Several programs connect information with spacetime geometry. Jacobson derived Einstein's field equations from the thermodynamics of local causal horizons~\cite{Jacobson1995}, Padmanabhan developed complementary thermodynamic formulations~\cite{Padmanabhan2010}, and Verlinde derived gravity from entropic forces on holographic screens~\cite{Verlinde2011}. The present framework shares their information-theoretic starting point, but it works with the classical probability laws of many agents, not with horizon or screen entropies, and it derives no gravitational dynamics. Tensor-network proposals read the geometry of the AdS/CFT correspondence off entanglement~\cite{Swingle2012}, and quantum-information approaches recover spatial geometry from the entanglement structure of Hilbert space~\cite{Cao2017}; the present construction shares their strategy of reading geometry off informational structure and implements it through classical information geometry and gauge transport. Constructor theory reformulates information in physical terms, characterizing information media by which transformations are possible, and rejects information as an a priori mathematical primitive~\cite{Deutsch2015}. The closest physics precedents for individual ingredients are the following. The information clock of Section~\ref{sec:time_info} is the natural-gradient analog of thermodynamic length, which bounds dissipation and is tied to Fisher information~\cite{SalamonBerry1983, Crooks2007}, and its distinguishability unit is Wootters' statistical distance~\cite{Wootters1981}. Caticha's proposals that space, and spacetime, carry the information geometry of a statistical manifold~\cite{Caticha2005, Caticha2019spacetime} anticipate the pullback of Section~\ref{sec:it_from_bit}, and the finding that an anti-de Sitter metric is a generic feature of Fisher geometries with symmetry~\cite{ErdmengerGrosvenorJefferson2020} cautions against over-reading the hyperbolic Gaussian metric. Quantum Darwinism explains the objectivity of classical records by their redundant proliferation among many observers~\cite{Zurek2009}; it is the quantum precedent for the consensus mechanism of Section~\ref{sec:collective_geometry}. A free-energy principle for generic quantum systems~\cite{Fields2022} is the closest published attempt at the level-zero extension listed as open in Section~\ref{sec:open_problems}.

A second line is gauge-theoretic free energy and Bayesian mechanics. Sengupta, Tozzi, Cooray, Douglas, and Friston recast the variational free-energy principle in gauge-theoretic language, with prediction errors acting as local forces and precision weighting related to gauge fields~\cite{Sengupta2016NeuronalGauge}. It is the closest precursor on this axis; the present construction differs in supplying explicit section fields $\phi_i \in \mathfrak{gl}(K_q)$, the inter-agent transport $\Omega_{ij}$, its action on covariances, and a multi-agent free energy coupled through relative entropies. Geometric Bayesian mechanics extends the free-energy principle to systems bounded by Markov blankets, sets of variables that screen internal from external states~\cite{friston2019particular, dacosta2021bayesian, ramstead2023bayesian, Sakthivadivel2022, KirchhoffParrPalaciosFristonKiverstein2018}, and it is the precedent for the multi-level blanket reading of agents in Section~\ref{sec:one_rung}. We adopt such blankets at every level and do not derive them; whether they exist across arbitrary scales is open, since the partition can depend on the coarse-graining. Bruineberg et al.\ distinguish an epistemic Pearl blanket from an ontological Friston blanket~\cite{bruineberg2022emperor}; the blanket here is indexed by the observer by design, which is the Pearl reading. Aguilera et al.\ show in linear stochastic systems that the Bayesian-mechanics reading holds only at weak solenoidal coupling~\cite{aguilera2022particular}, while the couplings $\beta_{ij}$ and $\gamma_{ij}$ here are dense and irreversible, so the belief coupling is motivated as a declared consensus objective (Section~\ref{sec:mixture_derivation}). Biehl et al.\ show that foundational arguments equivocate between blanket definitions~\cite{biehl2021technical}; we anchor the construction to one, the per-agent support $\chi_i$ on the domain $\mathcal{U}_i$, under which the statement that internal states parametrize beliefs is a modeling stipulation. Friston and Frith treat the synchronization of beliefs between communicating agents without a law of transport between their frames~\cite{FristonFrith2015duet}. Consensus coupled through edge maps on graphs is studied as the connection Laplacian~\cite{SingerWu2012} and as sheaf-Laplacian opinion dynamics~\cite{HansenGhrist2021Discourse}, and gauge-equivariant networks impose frame covariance on learned features~\cite{cohen2019gauge}; the present construction differs in coupling probability laws through a relative-entropy functional with state-dependent attention weights.

A third line takes the observer as primary. Relational quantum mechanics holds that quantum states and the values of variables are relative to other physical systems, any system serving as an observer, and describes only the information that systems have about each other~\cite{Rovelli1996}. With independent links on the interaction graph, network holonomy in the present construction obstructs a global frame, a belief-level counterpart of the question of how the perspectives of relational quantum mechanics are linked~\cite{AdlamRovelli2022}. QBism keeps outcomes personal and needs no consensus~\cite{Fuchs2014, DeBrotaFuchsSchack2020}; Fuchs names Wheeler's ``law without law'' as its inspiration and calls the resulting view participatory realism~\cite{Fuchs2017}. The present framework shares QBism's personalist starting point but not its quantum content, and it models Wheeler's third question, which QBism leaves to communication between agents, by consensus as a stationary point. Lahav and Neemeh propose a relativistic theory of consciousness in which phenomenal experience is frame-relative, and they write the transformation between cognitive frames for functionally identical systems~\cite{LahavNeemeh2022, LahavNeemeh2025}; the frame fields $\phi_i$ and transports $\Omega_{ij}$ supply an explicit map between frames of different agents at the level of belief transport (Section~\ref{sec:lahav_convergence}).

